% ============================================================
\chapter{Economic and Financial Models}
\label{ch:economic}
% ============================================================

\section{Introduction}

Economists have long envied physics for the precision of its models. In 1956, Robert Solow proposed an economic growth model of Newtonian elegance: capital accumulates, labour grows, and technology progresses, all governed by a differential equation whose steady state illuminates the wealth disparities between nations. This model earned him the Nobel Prize in 1987. But economics is also the domain of uncertainty: financial markets oscillate, option prices defy intuition. It is there that the Black--Scholes model, rooted in It\^o's stochastic calculus, revolutionized finance in 1973. This chapter weaves the link between deterministic and stochastic differential equations applied to economics and finance, from the Solow model to option pricing.

% ------------------------------------------------------------
\section{The Solow Growth Model}
\label{sec:solow}
% ------------------------------------------------------------

\begin{definition}[Solow model]
The Solow model describes long-run economic growth.  Let $K(t)$ be
capital, $L(t)$ labour, and $Y(t)$ output.  The production function
is Cobb--Douglas:
\[
  Y = AK^\alpha L^{1-\alpha}, \qquad 0<\alpha<1,
\]
where $A$ is total factor productivity (TFP).  Capital accumulates
according to:
\[
  \dot{K} = sY - \delta K,
\]
where $s\in(0,1)$ is the savings rate and $\delta>0$ the depreciation
rate.
\end{definition}

\subsection{Per-Capita Capital}

Assume $L(t)=L_0 e^{nt}$ (population growth at rate $n$).  Setting
$k=K/L$ (capital per worker) and $y=Y/L$:

\begin{theorem}[Fundamental Solow equation]
Capital per worker satisfies:
\[
  \dot{k} = sAk^\alpha - (n+\delta)k.
\]
\end{theorem}

\begin{proof}
$\dot{k} = \dot{K}/L - Kn/L = (sY-\delta K)/L - nk
= sy - \delta k - nk = sAk^\alpha - (n+\delta)k$.
\end{proof}

\begin{proposition}[Solow equilibrium]
The positive equilibrium $k^*$ satisfies $sAk^{*\alpha}=(n+\delta)k^*$:
\[
  k^* = \left(\frac{sA}{n+\delta}\right)^{1/(1-\alpha)}.
\]
This equilibrium is \textbf{globally stable} for $k>0$.
\end{proposition}

\begin{keyformulas}[title=\textbf{Solow equilibrium}]
\[
  k^* = \left(\frac{sA}{n+\delta}\right)^{1/(1-\alpha)}, \qquad
  y^* = A\left(\frac{sA}{n+\delta}\right)^{\alpha/(1-\alpha)}.
\]
\end{keyformulas}

\begin{codeblock}[title=\textbf{Solow model simulation}]
\begin{minted}{python}
import numpy as np
import matplotlib.pyplot as plt
from scipy.integrate import solve_ivp

A, alpha = 1.0, 0.3
s, delta, n = 0.2, 0.05, 0.02

def solow(t, k):
    return s * A * k**alpha - (n + delta) * k

k_star = (s * A / (n + delta))**(1 / (1 - alpha))
print(f"Equilibrium: k* = {k_star:.4f}")

fig, (ax1, ax2) = plt.subplots(1, 2, figsize=(13, 5))

for k0 in [0.5, 1, 2, 5, 10, 15]:
    sol = solve_ivp(solow, [0, 100], [k0],
                    t_eval=np.linspace(0, 100, 500))
    ax1.plot(sol.t, sol.y[0], label=f'$k_0={k0}$')

ax1.axhline(k_star, color='gray', ls='--', label=f'$k^*={k_star:.2f}$')
ax1.set_xlabel('$t$'); ax1.set_ylabel('$k(t)$')
ax1.legend(fontsize=8); ax1.set_title('Convergence to $k^*$')

k_vals = np.linspace(0, 15, 200)
ax2.plot(k_vals, s * A * k_vals**alpha, 'b-', label='$sAk^\\alpha$')
ax2.plot(k_vals, (n + delta) * k_vals, 'r-', label='$(n+\\delta)k$')
ax2.plot(k_star, (n+delta)*k_star, 'ko', markersize=8)
ax2.set_xlabel('$k$'); ax2.set_ylabel('Investment / Depreciation')
ax2.legend(); ax2.set_title('Solow Diagram')
plt.tight_layout(); plt.savefig('solow.pdf')
\end{minted}
\end{codeblock}

% ------------------------------------------------------------
\section{Supply-Demand Equilibrium}
\label{sec:supply_demand}
% ------------------------------------------------------------

\begin{definition}[Price adjustment model]
Let $D(p)$ be demand and $S(p)$ supply as functions of price $p$.
The price adjusts proportionally to excess demand:
\[
  \dot{p} = \kappa\bigl(D(p) - S(p)\bigr), \qquad \kappa > 0.
\]
\end{definition}

\begin{example}[Linear functions]
With $D(p)=a-bp$ and $S(p)=c+dp$ ($a,b,c,d>0$):
\[
  \dot{p} = \kappa\bigl((a-c)-(b+d)p\bigr).
\]
The equilibrium $p^*=(a-c)/(b+d)$ is globally stable.
\end{example}

\subsection{Cobweb Model}

\begin{definition}[Cobweb model]
In discrete time, producers set supply based on past price:
\[
  S_t = S(p_{t-1}), \qquad p_t = D^{-1}(S_t).
\]
\end{definition}

\begin{proposition}
If $\abs{S'(p^*)/D'(p^*)} < 1$, the model converges to equilibrium.
If $\abs{S'(p^*)/D'(p^*)} > 1$, oscillations diverge.
\end{proposition}

% ------------------------------------------------------------
\section{The Black--Scholes Model}
\label{sec:black_scholes}
% ------------------------------------------------------------

\begin{definition}[Geometric Brownian motion]
The price of a financial asset $S(t)$ follows geometric Brownian
motion:
\[
  dS = \mu S\,dt + \sigma S\,dW_t,
\]
where $\mu$ is the drift, $\sigma$ the volatility, and $W_t$ a
standard Brownian motion.
\end{definition}

\begin{theorem}[Black--Scholes equation]
The price $V(S,t)$ of a European option on an asset of price $S$
satisfies the PDE:
\[
  \frac{\partial V}{\partial t}
  + \frac{1}{2}\sigma^2 S^2\frac{\partial^2 V}{\partial S^2}
  + rS\frac{\partial V}{\partial S} - rV = 0,
\]
where $r$ is the risk-free rate.
\end{theorem}

\begin{keyformulas}[title=\textbf{Black--Scholes formula for a European call}]
\[
  C(S,t) = S\,\Phi(d_1) - Ke^{-r(T-t)}\Phi(d_2),
\]
where $K$ is the strike price, $T$ the maturity, $\Phi$ the standard
normal CDF, and
\[
  d_1 = \frac{\ln(S/K)+(r+\sigma^2/2)(T-t)}{\sigma\sqrt{T-t}},
  \qquad d_2 = d_1 - \sigma\sqrt{T-t}.
\]
\end{keyformulas}

\begin{codeblock}[title=\textbf{Black--Scholes formula in Python}]
\begin{minted}{python}
import numpy as np
from scipy.stats import norm
import matplotlib.pyplot as plt

def black_scholes_call(S, K, T, r, sigma):
    d1 = (np.log(S/K) + (r + sigma**2/2)*T) / (sigma*np.sqrt(T))
    d2 = d1 - sigma * np.sqrt(T)
    return S * norm.cdf(d1) - K * np.exp(-r*T) * norm.cdf(d2)

def black_scholes_put(S, K, T, r, sigma):
    d1 = (np.log(S/K) + (r + sigma**2/2)*T) / (sigma*np.sqrt(T))
    d2 = d1 - sigma * np.sqrt(T)
    return K * np.exp(-r*T) * norm.cdf(-d2) - S * norm.cdf(-d1)

S0, K, T, r, sigma = 100, 100, 1.0, 0.05, 0.2
print(f"Call = {black_scholes_call(S0, K, T, r, sigma):.4f}")
print(f"Put  = {black_scholes_put(S0, K, T, r, sigma):.4f}")

S_range = np.linspace(60, 140, 200)
calls = [black_scholes_call(S, K, T, r, sigma) for S in S_range]

fig, ax = plt.subplots(figsize=(8, 5))
ax.plot(S_range, calls, 'b-', lw=2, label='Call price')
ax.plot(S_range, np.maximum(S_range - K, 0), 'r--', label='Payoff')
ax.set_xlabel('$S$'); ax.set_ylabel('Call price')
ax.legend(); ax.set_title('European Call (Black-Scholes)')
plt.tight_layout(); plt.savefig('black_scholes.pdf')
\end{minted}
\end{codeblock}

% ------------------------------------------------------------
\section{Monte Carlo Simulation}
\label{sec:monte_carlo}
% ------------------------------------------------------------

\begin{definition}[Monte Carlo pricing]
Simulate $M$ price paths $S(T)$ and estimate the option price by
the discounted average payoff:
\[
  \hat{C} = e^{-rT}\frac{1}{M}\sum_{i=1}^M \max(S_i(T)-K,\;0).
\]
\end{definition}

\begin{codeblock}[title=\textbf{Monte Carlo for European call}]
\begin{minted}{python}
import numpy as np

S0, K, T, r, sigma = 100, 100, 1.0, 0.05, 0.2
M = 100000
np.random.seed(42)

Z = np.random.standard_normal(M)
S_T = S0 * np.exp((r - 0.5*sigma**2)*T + sigma*np.sqrt(T)*Z)
payoffs = np.maximum(S_T - K, 0)
C_mc = np.exp(-r*T) * np.mean(payoffs)
se = np.exp(-r*T) * np.std(payoffs) / np.sqrt(M)

print(f"MC price = {C_mc:.4f} +/- {1.96*se:.4f} (95% CI)")
print(f"BS price = {black_scholes_call(S0, K, T, r, sigma):.4f}")
\end{minted}
\end{codeblock}

% ------------------------------------------------------------
\section{Markowitz Portfolio Optimisation}
\label{sec:markowitz}
% ------------------------------------------------------------

\begin{definition}[Portfolio optimisation]
For $n$ assets with expected returns $\boldsymbol{\mu}\in\R^n$ and
covariance matrix $\Sigma$, the Markowitz optimal portfolio minimises
variance for a target return $r_c$:
\[
  \min_{\mathbf{w}} \;\mathbf{w}^\top\Sigma\mathbf{w}
  \quad\text{s.t.}\quad
  \mathbf{w}^\top\boldsymbol{\mu}=r_c,\quad
  \mathbf{w}^\top\mathbf{1}=1.
\]
\end{definition}

\begin{proposition}[Efficient frontier]
The set of optimal portfolios forms a parabola in the (risk, return)
plane, called the \textbf{efficient frontier}.
\end{proposition}

% ------------------------------------------------------------
\section{Leontief Input-Output Model}
% ------------------------------------------------------------

\begin{definition}[Leontief model]
Let $A=(a_{ij})$ be the matrix giving the amount of good $i$ needed
to produce one unit of good $j$.  Production equilibrium satisfies:
\[
  \mathbf{x} = A\mathbf{x} + \mathbf{d}
  \quad\Leftrightarrow\quad
  \mathbf{x} = (I-A)^{-1}\mathbf{d},
\]
where $\mathbf{d}$ is final demand.
\end{definition}

\begin{proposition}
If the spectral radius satisfies $\rho(A)<1$, then $(I-A)$ is
invertible and $(I-A)^{-1}=\sum_{k=0}^\infty A^k\geq 0$.
\end{proposition}

% ------------------------------------------------------------
\section{Limitations of Financial Models}
% ------------------------------------------------------------

\begin{warning}[title=\textbf{Limitations of Black--Scholes}]
\begin{itemize}
  \item Volatility $\sigma$ is not constant (volatility smile).
  \item Real returns have heavy tails (extreme events more frequent
    than predicted by the normal distribution).
  \item Markets are not always liquid and frictionless.
  \item Financial crises reveal nonlinear correlations that Gaussian
    models ignore.
\end{itemize}
\end{warning}

% ------------------------------------------------------------
\section{Exercises}
% ------------------------------------------------------------

\begin{exercise}
For the Solow model with $A=1$, $\alpha=0.3$, $s=0.25$, $\delta=0.05$,
$n=0.01$:
\begin{enumerate}
  \item Compute $k^*$ and $y^*$.
  \item Study the effect of increasing the savings rate $s$.
  \item Find the optimal savings rate (Phelps' golden rule).
\end{enumerate}
\end{exercise}

\begin{exercise}
Simulate the cobweb model with $D(p)=100-2p$ and $S(p)=-10+3p$.
Determine whether oscillations converge or diverge.
\end{exercise}

\begin{exercise}
Compute the ``Greeks'' of the European call (Delta, Gamma, Vega,
Theta, Rho) both analytically and numerically.
\end{exercise}

\begin{exercise}
Compare the Monte Carlo price of a European put with the
Black--Scholes formula for different numbers of paths $M$.  Plot
the error as a function of $M$ and verify $O(1/\sqrt{M})$ convergence.
\end{exercise}

\begin{exercise}
For a portfolio of 3 assets with given returns and covariances,
plot the Markowitz efficient frontier.
\end{exercise}
